\subsection{Tangent bundle}

8.1.1. Let \(r \in \mathbf{N}_{\mathbb{K}}\) and let X be a K-manifold of class \(\mathbf{C}^r\) . One denotes by T(X) the disjoint union of the tangent spaces \(\mathrm{T}_x(\mathrm{X})\) (5.5.1) for \(x \in \mathrm{X}\) , and \(\pi\) the canonical map from T(X) to X. Let \(c = (\mathrm{U}, \varphi, \mathrm{E})\) be a chart of the manifold X; for \(x \in \mathrm{U}\) and \(t \in \mathrm{T}_x(\mathrm{X})\) , set \(\zeta_c(x, t) = (x, \theta_c^{-1}(t))\) , the map \(\theta_c\) being that defined in 5.5.1. The triple \(c' = (\mathrm{U}, \zeta_c, \mathrm{E})\) is a vector chart (7.1.1) of T(X) (endowed with the map \(\pi: \mathrm{T}(\mathrm{X}) \to \mathrm{X}\) ). There exists on T(X) a structure of vector bundle with base X and of class \(\mathbf{C}^{r-1}\) , unique, such that, for every chart \(c\) of X, the corresponding vector chart \(c'\) is a vector chart of T(X) \(^{(1)}\) . Endowed with this structure, T(X) is called the tangent bundle of X. In particular, T(X) is endowed with a manifold structure of class \(\mathbf{C}^{r-1}\) , and the triple (T(X), X, \(\pi\) ) is a fibration (7.1.4).

Let U be an open subset of a Banach space E and let i be the canonical injection of U into E. The vector chart \(c' = (U, \zeta_c, E)\) associated with the chart \(c = (U, i, E)\) from U defines an isomorphism from T(U) onto the trivial bundle \(E_U\) (7.1.5); we identify these two bundles by means of this isomorphism.

If X is a pure manifold of type E (5.1.7), then T(X) is a pure manifold of type E \(\times\) E.

8.1.2. Let \(f: X \to Y\) be a morphism of manifolds of class \(\mathbf{C}^r\) . We define an \(f\)-morphism \(\mathbf{T}(f): \mathbf{T}(\mathbf{X}) \to \mathbf{T}(\mathbf{Y})\) of vector bundles of class \(\mathbf{C}^{r-1}\) by \(\mathbf{T}(f)(x, t) = (f(x), \mathbf{T}_x(f)(t))\) for \(x \in \mathbf{X}\) and \(t \in \mathbf{T}_x(\mathbf{X})\) .

If \(f = \operatorname{Id}_{\mathbf{X}}\) , one has \(\mathbf{T}(f) = \operatorname{Id}_{\mathbf{T}(\mathbf{X})}\) . For \(f: \mathbf{X} \to \mathbf{Y}\) and \(g: \mathbf{Y} \to \mathbf{Z}\) , one has \(\mathbf{T}(g \circ f) = \mathbf{T}(g) \circ \mathbf{T}(f)\) .

8.1.3. Let \(f: \mathbf{X} \to \mathbf{Y}\) be a morphism of manifolds of class \(\mathbf{C}^r\) , and let \(f': \mathbf{T}(\mathbf{X}) \to f^* \mathbf{T}(\mathbf{Y})\)

\(^{1}\) When K = R and r = 1, T(X) is a topological vector bundle (cf. Notations and conventions).

be the unique X-morphism such that \(\mathrm{T}(f) = g \circ f'\) , where \(g\) is the canonical morphism \(f^{*}\mathrm{T}(\mathrm{Y}) \to \mathrm{T}(\mathrm{Y})\) (cf. 7.2.4).

For \(f\) to be an immersion (resp. a submersion), it is necessary and sufficient that \(0 \to \mathrm{T}(\mathrm{X}) \xrightarrow{f'} f^* \mathrm{T}(\mathrm{Y})\) (resp. \(\mathrm{T}(\mathrm{X}) \xrightarrow{f'} f^* \mathrm{T}(\mathrm{Y}) \to 0\) ) be a locally direct exact sequence (7.5.6). If \(f\) is an immersion, the cokernel bundle of \(f'\) is called the normal or transverse bundle of \(f\) . If \(f\) is a submersion, the kernel bundle of \(f'\) is called the tangent bundle to the fibers of \(f\) , or the relative tangent bundle of X over Y, and is denoted \(\mathrm{T}(\mathrm{X}/\mathrm{Y})\) . Its fiber \(\mathrm{T}_x(\mathrm{X}/\mathrm{Y})\) at \(x \in \mathrm{X}\) is the tangent space at \(x\) to the submanifold \(f^{-1}(f(x))\) of X. The sequence

\[
0 \longrightarrow \mathrm{T} _ {x} (\mathrm{X} / \mathrm{Y}) \stackrel {i} {\longrightarrow} \mathrm{T} _ {x} (\mathrm{X}) \stackrel {\mathrm{T} _ {x} (f)} {\longrightarrow} \mathrm{T} _ {f (x)} (\mathrm{Y}) \longrightarrow 0
\]

where i is the canonical injection, is exact.

For \(f\) to be étale, it is necessary and sufficient that \(f'\) be an isomorphism.

When K is of characteristic 0, for f to be a subimmersion, it is necessary and sufficient that \(f'\) be locally direct.

8.1.4. Let \(X_{1}\) and \(X_{2}\) be two manifolds. We have \(\mathrm{T}(\mathrm{X}_{1} \times \mathrm{X}_{2}) = p_{1}^{*}\mathrm{T}(\mathrm{X}_{1}) \oplus p_{2}^{*}\mathrm{T}(\mathrm{X}_{2})\) , where \(p_{1}\) and \(p_{2}\) are the canonical projections (5.6.3).

8.1.5. Let \(n\) be an integer \(\geqslant 0\) ; the tangent bundle \(\mathrm{T}(\mathbf{K}^{n})\) of \(\mathbf{K}^{n}\) is identified with the trivial bundle with fiber \(\mathbf{K}^{n}\) (8.1.1). The constant sections defined by the elements of the canonical basis of \(\mathbf{K}^{n}\) form a frame of \(\mathrm{T}(\mathbf{K}^{n})\) .

If X is a manifold and \((u^{1},\ldots,u^{n})\) a coordinate system of X on an open set U of X (5.1.10), the preceding frame defines, by transport of structure, a frame of T(X) on U; it is called the tangent frame defined by \((u^{1},\ldots,u^{n})\) and is denoted \((\partial/\partial u^{1},\ldots,\partial/\partial u^{n})\) (cf. 5.5.8).

